\section{Homogeneous shape variations and local vacuum evolution}

\subsection{General matrix action and five-direction variation}

Eliminate torsion by the independent connection equation and use the GHY-completed LC action. In a fixed class-A coframe, $e=qe^{\mathsf B}\sigma$ has spatial curvature $f(n_{\mathsf B})/q^2$ and tracefree extrinsic curvature $\mathsf S$, with $D=N^{-1}d/dt$.

\stepcounter{equation}
\begin{equation}\tag{S-04}\label{eq:S-04}
\begin{aligned}
n_{\mathsf B}&=e^{\mathsf B}ne^{\mathsf B},\qquad
\mathsf S=\tfrac12e^{-\mathsf B}D(e^{2\mathsf B})e^{-\mathsf B},\\
L&=\frac{\mathcal V_{\Sigma}Nq^3}{2\kappa^2}
 \left[-6H^2+\mathrm{tr}(\mathsf S^2)+\frac{f(n_{\mathsf B})}{q^2}\right].
\end{aligned}
\end{equation}

The formula for $f$ is the spatial Ricci trace obtained from the Koszul connection in diagonal $n$, and extends by constant rotation to every real symmetric $n$. Counting both $\mathsf B,D\mathsf B$ as first order, expand the exponentials preserving noncommuting product order. Then $\mathsf S=D\mathsf B+O(3)$: quadratic commutator terms cancel between the two factors $e^{-\mathsf B}$. The potential is

\stepcounter{equation}
\begin{equation}\tag{S-05}\label{eq:S-05}
\begin{aligned}
f(n_{\mathsf B})&=6\chi-2\mathrm{tr}(X\mathsf B)+U_2+O(3),\\
U_2&=2[\mathrm{tr}(\mathsf Bn)]^2
 +2(\mathrm{tr}n)\mathrm{tr}(\mathsf B^2n)\\
&\hspace{1em}-4\mathrm{tr}(\mathsf Bn\mathsf Bn)
 -4\mathrm{tr}(\mathsf B^2n^2).
\end{aligned}
\end{equation}

Indeed, insert $n_{\mathsf B}=n+\mathsf Bn+n\mathsf B+\mathsf B^2n/2+n\mathsf B^2/2+\mathsf Bn\mathsf B+O(3)$ into $f$ to obtain each term. Tracefreeness gives (E-04). The five directions are $\mathsf B_{11}=\mathsf b_1,\mathsf B_{22}=\mathsf b_2,\mathsf B_{33}=-\mathsf b_1-\mathsf b_2$ and three independent off-diagonal components. Their velocity Gram form is $\dot{\mathsf b}_1^2+\dot{\mathsf b}_2^2+(\dot{\mathsf b}_1+\dot{\mathsf b}_2)^2+2(\dot{\mathsf b}_3^2+\dot{\mathsf b}_4^2+\dot{\mathsf b}_5^2)$, which is positive definite.

Diagonalize $n$ to classify $X=0$. The difference $X_i-X_j=(n_i-n_j)[2(n_i+n_j)-s_1]$ excludes three distinct eigenvalues. With two equal eigenvalues, write $n=(a,a,b)$; the difference is proportional to $b(b-a)$, so $b=a$ or $b=0$. Substitution into (E-04) verifies both branches, proving necessity and sufficiency in (E-05).

Independently evaluating the momentum equation about the common-scale background gives $\delta G_{01}=(n_2-n_3)D\mathsf B_{23}/q$ and cyclic permutations. General off-diagonal shape velocities are therefore not unrestricted. All momentum equations vanish for $n=0$, and are satisfied in the diagonal class-A sector with zero off-diagonal velocities.

\subsection{Torus lapse constraint and exact family}

For $n=0$ the potential vanishes identically. Substitute $q=q_0e^{\epsilon\rho}$, $\mathsf B=\epsilon\mathsf B_1$, $N=N_0+\epsilon\delta N$ into (S-04). The quadratic term is (S-01), while the cubic term and original lapse equation give

\stepcounter{equation}
\begin{equation}\tag{S-06}\label{eq:S-06}
\begin{aligned}
L^{[3]}&=\frac{A_{\Sigma}q_0^3}{6N_0}\left(3\rho-\frac{\delta N}{N_0}\right)
 [\mathrm{tr}(\dot{\mathsf B}_1^2)-6\dot\rho^2],\\
E_N^{[2]}&=\frac{A_{\Sigma}q_0^3}{6N_0^2}[6\dot\rho^2-\mathrm{tr}(\dot{\mathsf B}_1^2)]
 =\partial_{\delta N}L^{[3]}.
\end{aligned}
\end{equation}

Absence of the first-order lapse variation from the quadratic action does not remove the constraint. Vary the original lapse first, or vary the cubic action with respect to $\delta N$, to obtain the same second-order constraint. The quadratic Euler equations are $D^2\rho=D^2\mathsf B_1=0$, giving the affine solutions and cone in (S-01) as necessary conditions.

For sufficiency use (S-02). Before the constant congruence, mixed extrinsic curvature is $K=(3\epsilon h/\psi)\mathsf P$. Spatial Ricci vanishes and $DK+(\mathrm{tr}K)K=0$. The temporal Ricci component is

\stepcounter{equation}
\begin{equation}\tag{S-07}\label{eq:S-07}
\mathrm{Ric}_{00}=\frac{9\epsilon^2h^2}{\psi^2}[1-\mathrm{tr}(\mathsf P^2)]=-\frac{\epsilon^2}{\psi^2}[\mathrm{tr}(\Sigma^2)-6h^2].
\end{equation}

A constant congruence preserves local Ricci flatness, and every constant positive definite spatial metric descends to the same lattice. This does not declare a lattice transformation to be gauge. First-order matrix expansion gives $g_\epsilon=q_0^2[I+2\epsilon(\rho_0I+\mathsf B_0+h\tau I+\Sigma\tau)]+O(\epsilon^2)$, even if initial shape and $\Sigma$ do not commute. The determinant gives $q_\epsilon$ in (S-02). For $h=0$, positivity forces $\Sigma=0$, realized by a static family.

Any first-order lapse variation is realized by choosing positive $N_\epsilon=N_0+\epsilon\delta N+O(\epsilon^2)$ and composing the exact family with $\tau_\epsilon=\int N_\epsilon dt$. Since the background is static, this composition leaves its first-order spatial metric tangent unchanged. Local or initial-data comparisons differ from a two-endpoint fixed-metric variation problem; in the latter all affine scale and shape coefficients vanish.

\subsection{Diagonal closure and short-time coefficients}

Use diagonal $n$ and $e^i=qe^{\beta_i}\sigma^i$, $\sum_i\beta_i=0$. Put $m_i=n_ie^{2\beta_i}/q$ and evaluate $f,X_i$ below on physical $m$. The Koszul connection gives the diagonal spatial Ricci eigenvalues $2m_i^2-(\sum m)m_i+f$, whose tracefree parts are $X_i=2m_i^2-(\sum m)m_i+2f/3$. Extrinsic curvature is diagonal, and the momentum and off-diagonal Einstein equations vanish, so this sector closes.

\stepcounter{equation}
\begin{equation}\tag{S-08}\label{eq:S-08}
\begin{aligned}
Dq&=qH,& DH&=-3H^2-f/3,\\
D^2\beta_i&=-3HD\beta_i-X_i(m),& Dm_i&=m_i(2D\beta_i-H),\\
\mathcal C&=f+6H^2-\sum_i(D\beta_i)^2=0,&D\mathcal C&=-6H\mathcal C.
\end{aligned}
\end{equation}

Constraint propagation follows by differentiating (S-08) and using $Df=-2Hf-2\sum_iX_iD\beta_i$. On the constraint, $D^2q=-q\sum_i(D\beta_i)^2/3$. At $q=q_0,\beta=D\beta=0,H_0^2=-\chi/q_0^2$, one has $D^2\beta_i=-X_i(n)/q_0^2$ and $DX_i(m)=-2H_0X_i(n)/q_0^2$, hence $D^3\beta_i=5H_0X_i(n)/q_0^2$. Also $D^2q=D^3q=0$ and

\stepcounter{equation}
\begin{equation}\tag{S-09}\label{eq:S-09}
D^4q\big|_0=-\frac{2q_0}{3}\sum_i(D^2\beta_i|_0)^2=-\frac{2J_2(n)}{3q_0^3}.
\end{equation}

The supplement compares 12 points at $q_0=1$, Nil and Sol, both signs of $H_0$, and $\hat\tau=0.01,0.02,0.04$, using DOP853 with rtol $2\times10^{-13}$, atol $2\times10^{-15}$ and maximum step $0.002$. Constraint residuals and Taylor differences are saved. These point comparisons and local Taylor orders are distinct from a certified interval error bound.
